\documentclass[11pt]{article}

\usepackage[a4paper,margin=1in]{geometry}
\usepackage{amsmath,amssymb,bm}
\usepackage{enumitem}

\title{STL Dust Simulations for the BICEP/Keck Foreground Model Suite}
\author{}
\date{}

\begin{document}

\maketitle

\section{Scientific Context}

The BICEP/Keck program constrains the tensor-to-scalar ratio $r$ from
degree-scale CMB $B$-mode polarization.  At the sensitivity of BK18 and the
upcoming BK24 analysis, the dominant astrophysical systematic is Galactic
polarized foreground emission, especially dust.  The BICEP/Keck likelihood
fits multi-frequency $BB$ spectra with a model containing lensed $\Lambda$CDM,
tensor modes, dust, synchrotron, and noise.  In BK18, the baseline dust model
was a power-law angular spectrum and a modified-blackbody spectral energy
distribution (SED), with fitted dust amplitude
\[
  A_{d,353} \simeq 4.4\,\mu{\rm K}^2
\]
at $\ell=80$ and $353\,{\rm GHz}$ on the BICEP/Keck field.

The requested deliverable is not a deterministic reconstruction of the true
BICEP sky.  The requested deliverable is an ensemble of independent foreground
realizations whose statistics are realistic for the BICEP field and consistent
with the existing BICEP/Keck constraints.  These simulations will be passed
through the BICEP/Keck analysis machinery to test whether foreground complexity
not captured by the baseline parametric model can bias the recovered value of
$r$.

\section{Target Data Product}

The target product is an ensemble of synthesized foreground maps
\[
  \left\{
    \tilde{\bm s}_{\nu}^{\prime(a)}
  \right\}_{a=1}^{N_{\rm sim}},
  \qquad
  \tilde{\bm s}_{\nu}^{\prime(a)}
  =
  \left(\tilde I_\nu^{\prime(a)},
        \tilde Q_\nu^{\prime(a)},
        \tilde U_\nu^{\prime(a)}\right),
\]
on the BICEP/Keck sky patch, at the observing bandpasses supplied by
BICEP/Keck.  The desired ensemble size is
\[
  N_{\rm sim}\simeq 500,
\]
matching the standard BICEP/Keck simulation count.  The maps should be
provided at
\[
  N_{\rm side}=512,
\]
and should have controlled angular statistics over the approximate range
\[
  30 \lesssim \ell \lesssim 500.
\]
If necessary for a first delivery, the upper end may be reduced to
$\ell\simeq 300$, but the degree-scale range around $\ell\simeq 80$ is
essential.

The primary foreground component is polarized thermal dust.  Synchrotron may
also be included if a corresponding amplitude/SED prescription is agreed upon,
but the immediate problem is to generate dust realizations whose amplitude,
power spectrum, non-Gaussian morphology, and frequency decorrelation are
compatible with BICEP/Keck data.

\section{Input Data and Hypotheses}

Let $\Omega_{\rm HL}$ denote a collection of clean high-latitude sky patches
outside the Galactic plane.  The reference data are cleaned Planck polarized
dust maps at $353\,{\rm GHz}$,
\[
  \tilde{\bm s}_{353,p}
  =
  \left(\tilde Q_{353,p},\tilde U_{353,p}\right),
  \qquad p\in\Omega_{\rm HL},
\]
possibly together with an intensity dust tracer such as
\[
  \tilde I_{857,p}
\]
or another high signal-to-noise dust morphology template.  These are not
``training'' data in the machine-learning sense; they are the observed,
cleaned Planck dust patches from which we measure the target statistics.

The working hypotheses are:
\begin{enumerate}[label=(H\arabic*)]
  \item High-latitude Planck dust patches contain representative samples of
  the non-Gaussian polarized dust morphology relevant for the BICEP/Keck field.

  \item A finite set of scattering-transform statistics, supplemented by
  explicit power-spectrum constraints, captures the multiscale, anisotropic,
  non-Gaussian statistics needed for BICEP/Keck foreground validation.

  \item The generated realizations need not be phase-locked to the actual
  Planck sky at low multipoles.  In fact, BICEP/Keck prefers statistically
  independent realizations rather than maps deterministically tied to
  ``Planck truth'' on the BICEP field.  In other words, the goal is not to
  reproduce the particular observed sky realization pixel by pixel, but to
  reproduce the relevant ensemble statistics.

  \item The high-latitude maps must be amplitude-calibrated using a simple
  map-space normalization from the component-separated BICEP patch, and then
  spectrally extrapolated so that, on the BICEP/Keck mask, their amplitudes
  and cross-frequency decorrelation are compatible with the BK18 constraints.
  Cross-frequency decorrelation means that the dust pattern at two frequencies
  is not exactly the same map multiplied by two constants; spatial variations
  of the dust SED can suppress the cross-spectrum between frequencies relative
  to their auto-spectra.
\end{enumerate}

\section{STL Statistical Model at 353 GHz}

Let $\Phi$ denote the STL feature map.  In the current Planck component
separation code, $\Phi$ includes mean and variance, scattering covariance
coefficients
\[
  S_1,\ S_2,\ S_3,\ S_4,
\]
and, when enabled, radially binned Fourier power spectra.  The statistics may
include selected cross-statistics between channels, for example between
polarized maps and an auxiliary intensity tracer.

The generative problem at $353\,{\rm GHz}$ can be written as follows.  We seek
a probability law $\mathcal P_{\rm dust}$ on synthesized BICEP-patch
polarization maps
\[
  \tilde{\bm s}_{353}^{\prime(a)}
  =
  \left(\tilde Q_{353}^{\prime(a)},\tilde U_{353}^{\prime(a)}\right)
\]
such that its expected STL statistics match those measured from the
high-latitude Planck reference data:
\[
  \mathbb E_{\tilde{\bm s}'\sim\mathcal P_{\rm dust}}
  \left[\Phi(\tilde{\bm s}', \tilde I_{\rm aux}')\right]
  \simeq
  \frac{1}{N_{\rm HL}}
  \sum_{p\in\Omega_{\rm HL}}
  \Phi(\tilde{\bm s}_{353,p}, \tilde I_{{\rm aux},p}).
\]
The generated maps are then checked against the BICEP/Keck two-point
constraints,
\[
  D_\ell^{BB,353}
  =
  \frac{\ell(\ell+1)}{2\pi} C_\ell^{BB,353}
  \simeq
  A_{d,353}
  \left(\frac{\ell}{80}\right)^{\alpha_d},
\]
but this is not intended to be a hard per-$\ell$ constraint in the synthesis
objective.  The primary calibration is instead a simple map-space amplitude
renormalization, using the component-separated BICEP patch as the target
amplitude reference.

Operationally, one may generate maps by solving, for each realization $a$, an
optimization problem of the form
\[
  \tilde{\bm s}_{353}^{\prime(a)}
  =
  \operatorname*{arg\,min}_{\bm u}
  \left\|
    \Phi(\bm u,\tilde I_{\rm aux}^{\prime(a)})
    -
    \widehat{\Phi}_{\rm HL}
  \right\|_2^2,
\]
where $\widehat{\Phi}_{\rm HL}$ is the empirical STL summary of the
high-latitude Planck reference data, after the chosen map-space amplitude
normalization.  Random initial conditions, random reference patch draws, or
minibatch statistics provide independent realizations.

\section{Frequency Scaling and Decorrelation}

BICEP/Keck does not only need $353\,{\rm GHz}$ maps.  It needs maps integrated
over its observing bandpasses.  A minimal dust SED model is a modified
blackbody:
\[
  \tilde Q_\nu'(\hat n)
  =
  \tilde Q_{353}'(\hat n)\,
  f_d(\nu;\beta_d,T_d),
  \qquad
  \tilde U_\nu'(\hat n)
  =
  \tilde U_{353}'(\hat n)\,
  f_d(\nu;\beta_d,T_d),
\]
where
\[
  f_d(\nu;\beta_d,T_d)
  =
  \frac{g_{\rm CMB}(353)}{g_{\rm CMB}(\nu)}
  \left(\frac{\nu}{353\,{\rm GHz}}\right)^{\beta_d}
  \frac{B_\nu(T_d)}{B_{353}(T_d)}.
\]
This single-SED model produces no cross-frequency decorrelation: each
frequency map has exactly the same morphology, only rescaled.  A minimal
decorrelating example is
\[
  \tilde Q_\nu'(\hat n)
  =
  \tilde Q_{353}'(\hat n)
  \left(\frac{\nu}{353\,{\rm GHz}}\right)^{\beta_d(\hat n)}.
\]
If $\beta_d(\hat n)$ varies across the sky, two pixels with the same
$353\,{\rm GHz}$ amplitude may be rescaled differently at $150\,{\rm GHz}$.
The $150\,{\rm GHz}$ and $353\,{\rm GHz}$ maps are then no longer identical
up to a global multiplicative factor.


The email exchange
explicitly asks for enough SED complexity to reproduce the expected
decorrelation constraints, so the proposed model should include at least one
of the following extensions:
\begin{enumerate}[label=(S\arabic*)]
  \item spatially varying spectral index and/or temperature,
  $\beta_d(\hat n)$ and $T_d(\hat n)$;

  \item a multi-layer dust model,
  \[
    \tilde Q_\nu'(\hat n)
    =
    \sum_{k=1}^{K}
    \tilde Q_{k,353}'(\hat n)\,
    f_d(\nu;\beta_{d,k}(\hat n),T_{d,k}(\hat n)),
  \]
  with the same expression for $\tilde U_\nu'$;

  \item stochastic small-scale SED fluctuations whose amplitude is tuned to
  produce an allowed level of decorrelation.
\end{enumerate}

For two frequencies $\nu_1,\nu_2$, define the dust decorrelation coefficient
on the BICEP/Keck mask by
\[
  \Delta_d(\nu_1,\nu_2;\ell)
  =
  \frac{
    C_\ell^{BB,d}(\nu_1\times\nu_2)
  }{
    \sqrt{
      C_\ell^{BB,d}(\nu_1\times\nu_1)
      C_\ell^{BB,d}(\nu_2\times\nu_2)
    }
  }.
\]
BK18 found that the likelihood for the decorrelation parameter peaks at unity,
i.e. at no detected decorrelation, and Fig.~18 of BK18 provides the relevant
constraint for the $217\times353\,{\rm GHz}$ comparison.  Therefore the SED
model should be rich enough to explore decorrelation, but tuned so that
\[
  \Delta_d
\]
remains compatible with the BICEP/Keck likelihood on the sky mask.  Strongly
decorrelated models are useful stress tests only if clearly labeled as such.

\section{Ratio-Map SED Generative Strategy}

Following the discussion with Fran\c{c}ois Boulanger, Fran\c{c}ois Levrier,
and S\'ebastien, the preferred practical route is to separate the problem into
two coupled generative pieces:
\begin{enumerate}[label=(R\arabic*)]
  \item a generative model for the reference dust morphology, for instance at
  $353\,{\rm GHz}$;

  \item a generative model for frequency-ratio maps, which encode the spatial
  variation of the dust SED and therefore the possible decorrelation between
  channels.
\end{enumerate}
The ratio model is intended to be learned from simulations of pairs of
frequency channels.  Direct ratios of $Q$ or $U$ are undesirable because these
fields change sign and have many zero crossings.  We therefore define the
ratio on polarized intensity,
\[
  \tilde P_\nu(\hat n)
  =
  \left[
    \tilde Q_\nu^2(\hat n)
    +
    \tilde U_\nu^2(\hat n)
  \right]^{1/2},
\]
possibly after smoothing, masking, or debiasing if required.  For two
channels $353$ and $217\,{\rm GHz}$, the basic ratio field is
\[
  \rho_{353/217}(\hat n)
  =
  \mathcal R_\epsilon
  \left[
    \tilde P_{353}(\hat n),
    \tilde P_{217}(\hat n)
  \right],
\]
where $\mathcal R_\epsilon$ denotes the chosen regularized ratio operation.
In the simplest implementation,
\[
  \rho_{353/217}(\hat n)
  =
  \frac{
    \tilde P_{353}(\hat n)
  }{
    \tilde P_{217}(\hat n)+\epsilon_P
  },
\]
with a small regularization or mask to avoid unstable ratios in low-signal
pixels.  An equivalent log-ratio parameterization,
\[
  \eta_{353/217}(\hat n)
  =
  \log \rho_{353/217}(\hat n),
\]
may be preferable numerically because it keeps the generated ratio positive.
The essential object is a spatial field that describes how the polarized dust
amplitude changes between two frequencies.

Let $\Psi$ denote the scattering feature map used for ratio fields.  From an
ensemble of simulated ratios one measures
\[
  \widehat\Psi_{353/217}
  =
  \frac{1}{N_{\rm sim}^{\rm ratio}}
  \sum_{b=1}^{N_{\rm sim}^{\rm ratio}}
  \Psi\!\left(\rho_{353/217}^{(b)}\right).
\]
The ratio realizations used in the BICEP simulations are then generated by
solving
\[
  \rho_{353/217}^{\prime(a)}
  =
  \operatorname*{arg\,min}_{v}
  \left\|
    \Psi(v)
    -
    \widehat\Psi_{353/217}
  \right\|_2^2
  +
  \lambda_{\rm ratio}
  \,\mathcal C_{\rm ratio}(v),
\]
where $\mathcal C_{\rm ratio}$ denotes any additional constraints imposed on
the ratio maps, such as smoothness, allowed dynamic range, or a target
decorrelation level.

The $353\,{\rm GHz}$ morphology is still generated from the high-latitude dust
statistics.  Let
\[
  \mu_{\rm HL}
  =
  \frac{1}{N_{\rm HL}}
  \sum_{p\in\Omega_{\rm HL}}
  \Phi(\tilde{\bm s}_{353,p},\tilde I_{{\rm aux},p})
\]
be the mean scattering summary over high-latitude sky patches.  Before using
this target for BICEP/Keck simulations, the high-latitude maps are multiplied
by a single scalar chosen from the component-separated BICEP patch.  For
example, define a robust polarized-intensity amplitude
\[
  A_P[\tilde{\bm s}]
  =
  {\rm RMS}_{W_{\rm BK}}
  \left[
    \left(\tilde Q^2+\tilde U^2\right)^{1/2}
  \right],
\]
with the same smoothing, masking, and apodization convention applied to the
high-latitude patches and to the BICEP patch.  If
$\tilde{\bm s}_{353}^{\rm BK,cs}$ denotes the component-separated BICEP-patch
dust solution, set
\[
  \alpha_{\rm HL\to BK}
  =
  \frac{
    A_P[\tilde{\bm s}_{353}^{\rm BK,cs}]
  }{
    A_P[\tilde{\bm s}_{353}^{\rm HL}]
  },
  \qquad
  \tilde{\bm s}_{353,p}
  \leftarrow
  \alpha_{\rm HL\to BK}\,
  \tilde{\bm s}_{353,p}.
\]
The BICEP-adapted scattering target is then
\[
  \mu_{\rm HL}^{\rm BK}
  =
  \frac{1}{N_{\rm HL}}
  \sum_{p\in\Omega_{\rm HL}}
  \Phi(\alpha_{\rm HL\to BK}\tilde{\bm s}_{353,p},
       \alpha_{\rm HL\to BK}\tilde I_{{\rm aux},p}),
\]
or the corresponding normalized version if the auxiliary channel is not to be
rescaled.  The reference maps are generated as
\[
  \tilde{\bm s}_{353}^{\prime(a)}
  =
  \operatorname*{arg\,min}_{\bm u}
  \left\|
    \Phi(\bm u,\tilde I_{\rm aux}^{\prime(a)})
    -
    \mu_{\rm HL}^{\rm BK}
  \right\|_2^2,
\]
where $\mu_{\rm HL}^{\rm BK}$ denotes the high-latitude scattering target
renormalized to the BICEP/Keck amplitude.  There is no direct $BB$-spectrum
fitting term in this first implementation.  For an even simpler first
implementation, one does not necessarily need to generate the reference
morphology from scratch: one can start directly from the already solved
component-separation result on the BICEP patch itself and use it as the
initial $353\,{\rm GHz}$ dust template to test the extrapolation pipeline.

Given a synthesized reference map and a synthesized ratio map, one obtains the
lower-frequency polarized amplitude by inversion of the ratio:
\[
  \tilde P_{217}^{\prime(a)}(\hat n)
  =
  \frac{
    \tilde P_{353}^{\prime(a)}(\hat n)
  }{
    \rho_{353/217}^{\prime(a)}(\hat n)
  }.
\]
The simplest extrapolation keeps the polarization angle of the reference map
and rescales the polarized vector:
\[
  \begin{pmatrix}
    \tilde Q_{217}^{\prime(a)}(\hat n) \\
    \tilde U_{217}^{\prime(a)}(\hat n)
  \end{pmatrix}
  =
  \frac{
    \tilde P_{217}^{\prime(a)}(\hat n)
  }{
    \tilde P_{353}^{\prime(a)}(\hat n)+\epsilon_P
  }
  \begin{pmatrix}
    \tilde Q_{353}^{\prime(a)}(\hat n) \\
    \tilde U_{353}^{\prime(a)}(\hat n)
  \end{pmatrix}.
\]
More generally, for any target band $\nu$ one can learn or prescribe a ratio
field $\rho_{353/\nu}$ and set
\[
  \tilde P_{\nu}^{\prime(a)}(\hat n)
  =
  \frac{
    \tilde P_{353}^{\prime(a)}(\hat n)
  }{
    \rho_{353/\nu}^{\prime(a)}(\hat n)
  },
\]
with the same polarized-vector rescaling used to obtain
$\tilde Q_\nu^{\prime(a)}$ and $\tilde U_\nu^{\prime(a)}$.  This makes the SED
extrapolation a generative problem for polarized-intensity ratio fields rather
than only a parametric modified-blackbody scaling.  The ratio statistics
control how much the polarized amplitude changes between channels, and
therefore contribute to the decorrelation coefficients measured on the
BICEP/Keck mask.

The validation requirement is unchanged: after extrapolation to the relevant
BICEP/Keck bandpasses, the maps must satisfy the amplitude, angular-spectrum,
and decorrelation constraints.  In particular, the ratio-generated maps should
be checked through
\[
  \Delta_d(\nu_1,\nu_2;\ell)
  =
  \frac{
    C_\ell^{BB,d}(\nu_1\times\nu_2)
  }{
    \sqrt{
      C_\ell^{BB,d}(\nu_1\times\nu_1)
      C_\ell^{BB,d}(\nu_2\times\nu_2)
    }
  }
\]
on the BICEP/Keck mask and compared with the BK18/BK24 allowed range.

\subsection{Working in $Q,U$ while constraining $BB$}

The generative variables in the proposed pipeline are polarization maps,
\[
  \tilde{\bm s}_{\nu}^{\prime(a)}
  =
  \left(
    \tilde Q_{\nu}^{\prime(a)},
    \tilde U_{\nu}^{\prime(a)}
  \right),
\]
or quantities derived from them such as the polarized intensity
$\tilde P_\nu'=(\tilde Q_\nu'^2+\tilde U_\nu'^2)^{1/2}$.  The BICEP/Keck
foreground constraints, however, are naturally stated in terms of $E/B$
spectra, in particular the dust $BB$ amplitude at $\ell\simeq80$ and the
shape of $D_\ell^{BB}$ over the BICEP/Keck range.

The intended strategy is therefore not to microtune the $BB$ spectrum mode by
mode, nor to use $BB$ as the primary amplitude calibration.  The primary
amplitude calibration is the single map-space factor
$\alpha_{\rm HL\to BK}$ defined from the component-separated BICEP patch.
After this calibration, each synthesized $Q,U$ realization is transformed to
$E,B$ on the BICEP/Keck patch, with the same mask/apodization convention used
for the validation spectra, and its $BB$ spectrum is measured:
\[
  \widehat D_\ell^{BB,(a)}
  =
  \widehat D_\ell^{BB}
  \left[
    W_{\rm BK}
    \left(
      \tilde Q_{\nu}^{\prime(a)},
      \tilde U_{\nu}^{\prime(a)}
    \right)
  \right].
\]
This measured spectrum is then used as a diagnostic only.  It verifies that
the map-space calibration gives a $BB$ amplitude and slope compatible with
BICEP/Keck, but it is not used to fit or overwrite the realization.

This distinction is important.  The scattering model should determine the
polarized morphology in $Q,U$, and the resulting $BB$ spectrum should emerge
from that morphology after the usual $Q/U\rightarrow E/B$ projection.  The
$BB$ constraint is used to flag or reject realizations whose spectra are
grossly incompatible with BICEP/Keck data.  It should not be used as a hard
per-$\ell$ target that overwrites the non-Gaussian structure or
frequency-dependent morphology generated by the model.

\section{Calibration and Validation Against BICEP/Keck Constraints}

The operational calibration is deliberately simple.  We first use the
component-separated BICEP patch to define a single map-space amplitude factor,
and we multiply the high-latitude $Q,U$ maps by that number before measuring
the BICEP-adapted scattering target.  The $BB$ spectrum is then computed only
after synthesis and SED extrapolation, as a validation product.

At minimum, on the BICEP/Keck mask:
\begin{enumerate}[label=(C\arabic*)]
  \item the map-space polarized-intensity amplitude at $353\,{\rm GHz}$ should
  match the component-separated BICEP-patch reference, through the single
  scalar $\alpha_{\rm HL\to BK}$;

  \item after this calibration, the resulting $BB$ spectrum should be checked
  against the BK18 scale, approximately
  $A_{d,353}\simeq 4.4\,\mu{\rm K}^2$ at $\ell=80$, or the updated BK24 target
  if supplied;

  \item the angular spectrum should have an acceptable slope over
  $30\lesssim\ell\lesssim500$, but it is not microtuned bin by bin;

  \item the cross-frequency decorrelation coefficient should lie inside the
  allowed BICEP/Keck range;

  \item the maps should not be locked to the observed BICEP/Planck phases at
  low multipoles;

  \item the ensemble scatter across realizations should be available, since
  BICEP/Keck will use the simulations to estimate bias and uncertainty in
  recovered $r$.
\end{enumerate}

\section{Mathematical Workflow and Problem Statement}

The proposed workflow can be summarized as a sequence of maps and constraints.
First, start from cleaned high-latitude Planck polarization patches
\[
  \left\{
    \tilde{\bm s}_{353,p}
    =
    (\tilde Q_{353,p},\tilde U_{353,p})
  \right\}_{p\in\Omega_{\rm HL}},
  \qquad
  \tilde I_{{\rm aux},p}.
\]
Measure their empirical STL summary
\[
  \widehat{\Phi}_{\rm HL}
  =
  \frac{1}{N_{\rm HL}}
  \sum_{p\in\Omega_{\rm HL}}
  \Phi(\tilde{\bm s}_{353,p},\tilde I_{{\rm aux},p}),
\]
after applying the single map-space normalization.  The normalization factor
$\alpha_{\rm HL\to BK}$ is inferred from the component-separated BICEP patch.

Second, synthesize independent BICEP-patch dust realizations at
$353\,{\rm GHz}$:
\[
  \tilde{\bm s}_{353}^{\prime(a)}
  =
  \operatorname*{arg\,min}_{\bm u}
  \left[
    \left\|
      \Phi(\bm u,\tilde I_{\rm aux}^{\prime(a)})
      -
      \widehat{\Phi}_{\rm HL}
    \right\|_2^2
  \right].
\]
The prime distinguishes synthesized maps from the cleaned Planck reference
maps.  The index $a$ labels independent statistical realizations.

Third, validate the amplitude of the synthesized realization on the
BICEP/Keck patch.  If $W_{\rm BK}(\hat n)$ is the BICEP/Keck mask or
apodization window, measure both the map-space polarized-intensity amplitude
and the induced $BB$ amplitude:
\[
  \widehat A_{P,353}^{(a)}
  =
  A_P
  \left[
    W_{\rm BK}\tilde{\bm s}_{353}^{\prime(a)}
  \right],
  \qquad
  \widehat D_\ell^{BB,(a)}
  =
  \widehat D_\ell^{BB}
  \left[
    W_{\rm BK}\tilde{\bm s}_{353}^{\prime(a)}
  \right].
\]
No $BB$-based rescaling is applied at this stage in the baseline workflow.
The $BB$ spectrum is kept as a validation diagnostic for compatibility with
BICEP/Keck.

Fourth, apply an SED model to obtain bandpass maps
\[
  \left\{
    \tilde{\bm s}_{\nu}^{\prime(a)}
  \right\}_{\nu\in\mathcal B},
\]
where $\mathcal B$ is the set of BICEP/Keck observing bandpasses.  The SED
parameters must be tuned so that the decorrelation coefficients
\[
  \Delta_d^{(a)}(\nu_1,\nu_2;\ell)
  =
  \frac{
    C_\ell^{BB,d}
    \!\left[
      W_{\rm BK}\tilde{\bm s}_{\nu_1}^{\prime(a)}
      \times
      W_{\rm BK}\tilde{\bm s}_{\nu_2}^{\prime(a)}
    \right]
  }{
    \sqrt{
      C_\ell^{BB,d}
      \!\left[
        W_{\rm BK}\tilde{\bm s}_{\nu_1}^{\prime(a)}
        \times
        W_{\rm BK}\tilde{\bm s}_{\nu_1}^{\prime(a)}
      \right]
      C_\ell^{BB,d}
      \!\left[
        W_{\rm BK}\tilde{\bm s}_{\nu_2}^{\prime(a)}
        \times
        W_{\rm BK}\tilde{\bm s}_{\nu_2}^{\prime(a)}
      \right]
    }
  }
\]
remain compatible with the BK18/BK24 allowed range.

In compact form, the mathematical problem is therefore to construct an
ensemble distribution $\mathcal P_{\Theta}$ for
\[
  \left\{
    \tilde{\bm s}_{\nu}^{\prime}
  \right\}_{\nu\in\mathcal B}
\]
such that, on the BICEP/Keck patch,
\[
  \mathbb E_{\mathcal P_\Theta}
  \left[
    \Phi(\tilde{\bm s}_{353}^{\prime},\tilde I_{\rm aux}^{\prime})
  \right]
  \approx
  \widehat{\Phi}_{\rm HL},
\]
\[
  A_P[\tilde{\bm s}_{353}^{\prime}]
  \approx
  A_P[\tilde{\bm s}_{353}^{\rm BK,cs}],
\]
with the amplitude matching imposed through the single map-space factor
$\alpha_{\rm HL\to BK}$.  The resulting $BB$ spectrum is then checked a
posteriori,
\[
  D_\ell^{BB,d}
  \ \text{compatible with}\
  A_{d,353}^{\rm BK}
  \left(\frac{\ell}{80}\right)^{\alpha_d},
  \qquad
  30\lesssim \ell \lesssim 500,
\]
but it is not fitted directly in the synthesis objective.  Finally,
\[
  \Delta_d(\nu_1,\nu_2;\ell)
  \in
  \mathcal C_{\rm BK},
\]
where $\mathcal C_{\rm BK}$ denotes the BICEP/Keck decorrelation constraint.
The scientific test is then to run the BICEP/Keck likelihood on
\[
  {\rm CMB}^{(a)} + {\rm noise}^{(a)}
  + \tilde{\bm s}_{\nu}^{\prime(a)}
\]
and measure the induced shift
\[
  \Delta r
  =
  \widehat r_{\rm recovered}
  -
  r_{\rm input}.
\]
The model is successful if it provides foreground realizations that are
realistic, independent, BICEP-compatible, and useful for quantifying whether
non-Gaussian dust morphology or SED complexity can bias $r$ at a relevant
level.

\end{document}
